\section{Related Work and Positioning}
\label{sec:related}

This review sits at the intersection of three literatures that have, until now, developed largely in isolation: legal and jurisdictional analyses of AI regulation, technical surveys of AI evaluation methods, and emerging efforts to standardize transparency and documentation. We position our contribution against each in turn and then state the specific gap our methods-to-requirements mapping is designed to close.

\subsection{Legal Overviews and Framework Comparisons}

A first body of work characterizes the global regulatory landscape and compares governance instruments at a conceptual level. The \textit{Global AI Governance Overview} provides a broad synthesis of regulatory requirements across jurisdictions, cataloguing obligations and contrasting the risk-tiered, rights-based posture of the EU AI Act \cite{euaiact2024} with the voluntary, risk-management orientation of the NIST AI RMF \cite{nist2023airmf} and the management-system structure of ISO/IEC 42001 \cite{iso42001}, among others \cite{global2024governance}. Complementary cross-regional studies analyze how risk management frameworks differ across the EU, U.S., UK, and China, emphasizing institutional and policy divergences \cite{almaamari2025crossregional}. These overviews are valuable for situating obligations, but they remain at the level of legal text and policy intent; they do not descend to the question of \textit{which technical procedures} an engineer would execute to demonstrate that a system satisfies a given clause.

Closer to our aim is structured mapping work that connects regulatory concepts to standards artifacts. The open knowledge-graph approach of Hernandez et al. links concepts and requirements between the EU AI Act and international standards, enabling machine-readable traceability between legal provisions and harmonized standards \cite{hernandez2024knowledgegraph}. Related efforts develop ontologies for fundamental-rights impact assessments \cite{rintamaki2025ontology} and examine how fundamental rights are operationalized through European standardization \cite{hodac2024fundamentalrights, meding2025algorithmicdiscrimination}. This line of research maps \textit{requirements to standards and concepts}; our work is orthogonal and complementary, mapping \textit{requirements to executable technical evaluation methods}. Where a knowledge graph tells a practitioner which standard governs a clause, our matrix tells them which test, benchmark, or measurement procedure produces the evidence that clause demands.

\subsection{Per-Domain Evaluation Surveys}

A second, substantial literature surveys the technical methods themselves. General LLM-evaluation surveys catalogue capability, robustness, and task benchmarks \cite{chang2024survey, liang2023holistic}, while trustworthiness-focused studies assemble dimensions such as safety, fairness, privacy, and truthfulness and provide accompanying assessments \cite{huang2024trustllm, wang2023decodingtrust}. A rapidly growing strand surveys the evaluation of LLM-based and agentic systems, addressing autonomy, tool use, and multi-agent interaction \cite{yehudai2025survey, mohammadi2025evaluation}, alongside reviews of trust, risk, and security threats and countermeasures for agents \cite{yu2025survey, gan2024navigating, raza2025trism, qi2026trustworthy}. These surveys are thorough in enumerating and taxonomizing \textit{methods}, and broader responsible-AI surveys situate them among frameworks and best practices \cite{gadekallu2025responsibleaisurvey}. Crucially, however, they organize methods by technical property or threat model rather than by regulatory obligation: they answer \textit{what can be measured} but leave implicit \textit{which measurement discharges which legal requirement}. The mapping from a fairness benchmark or a robustness probe to a specific Article of the EU AI Act, a NIST RMF function, or an ISO/IEC 42001 control is left to the reader.

\subsection{Transparency and Documentation Efforts}

A third strand pursues transparency and structured documentation as a governance mechanism. The Foundation Model Transparency Index scores developer disclosures across dimensions spanning data, compute, and downstream impact \cite{wan2025fmti, raffel2024transparency}, and audit-card and model-card proposals seek to contextualize and standardize how evaluations and their provenance are reported \cite{staufer2025auditcards, puhlfurss2025modelcards}. Parallel work supplies technical-documentation templates and continuous-assurance postures aligned to the AI Act \cite{lucaj2025techops, lupo2026trustworthyposture}. These efforts standardize the \textit{reporting} of evaluations and improve their comparability, but they presuppose that the underlying technical methods and their regulatory relevance have already been selected; they describe how to document evidence rather than how to choose the method that generates it.

\subsection{The Gap We Fill}

Taken together, the prior literature offers legal overviews that stop above the technical layer, evaluation surveys that enumerate methods without binding them to obligations, and documentation schemas that report results without prescribing methods. No existing work, to our knowledge and within the technical and engineering literature as of our search window through mid-2026, provides a rigorous mapping from concrete AI-evaluation methods to specific regulatory requirements across the NIST AI RMF \cite{nist2023airmf}, the EU AI Act \cite{euaiact2024}, and ISO/IEC 42001 \cite{iso42001}; legal and policy scholarship has examined the obligations without this technical-methods focus. This review supplies that mapping, building on prior, unpublished structured evaluation frameworks for high-stakes and agentic systems (citations withheld for double-blind review), and consolidates it in the methods $\times$ requirements matrix of Table~\ref{tab:matrix}. Table~\ref{tab:diff} summarizes how our contribution differs from the most closely related prior work along the axes of regulatory coverage, technical granularity, and requirement-level mapping.
